\chapter{PENDAHULUAN}\label{pendahuluan}
\section{Latar Belakang}\label{latar belakang}

Berisi uraian tentang gambaran permasalahan/kebutuhan dan hal-hal yang mendasari pentingnya dilakukan penelitian ini, atau alasan mengapa penelitian ini penting untuk dilakukan.

\section{Perumusan Masalah}
Berisi uraian tentang:
\begin{enumerate}
    \item identifikasi akar permasalahan/kebutuhan, dan 
   \item pendekatan (\textit{approach}) penyelesaiannya.
\end{enumerate}


Identifikasi akar permasalahan/kebutuhan perlu dilakukan karena apa yang tampak di permukaan sebagai permasalahan/kebutuhan belum tentu merupakan akar permasalahan/kebutuhan yang sesungguhnya.  Ia bisa jadi muncul sebagai efek dari sebab yang lebih mendasar. Padahal, penyelesaian terhadap masalah/kebutuhan yang dirumuskan secara tidak tepat tidak akan menyelesaikan masalah atau menjawab kebutuhan. Oleh karena itu, agar diperoleh penyelesaian efektif, maka masalah/kebutuhan harus dirumuskan dengan tepat.

Pendekatan penyelesaian masalah menggambarkan secara singkat cara atau metode yang akan ditempuh untuk menyelesaikan akar permasalahan yang telah berhasil diidentifikasi. Cakupan permasalahan bisa sangat luas, karena ia bisa dilihat dari berbagai sudut pandang, dan boleh jadi tidak bisa dilakukan pelaksanaannya dalam kerangka Tugas Akhir. Oleh karena itu, pendekatan yang diusulkan perlu dibatasi  menurut cara pandang tertentu yang dianggap memadai atau layak.

\section{Tujuan}
Berisi uraian tentang tujuan yang akan dicapai dalam penelitian.

\section{Manfaat}
Berisi uraian tentang manfaat yang dapat diperoleh bila  tujuan penelitian tercapai.

